% from ega2-1-fr.tex
\oldpage[II]{68}

\begin{proof}
La seconde assertion est un cas particulier de
(\hyperref[II.3.7.6-fr]{3.7.6})~; d'autre part, pour montrer que
$r_{\mathcal{L},\psi}$ est dominant, on peut se borner au cas où $Y$ est
affine~; si $s\in S$ est telle que $\psi^\flat(s)$ soit localement
nilpotente, comme on peut écrire $s=\varphi_\lambda(s_\lambda)$ pour un
$\lambda$ au moins, on conclut de l'hypothèse et de
(\hyperref[II.3.7.5-fr]{3.7.5}) que $s_\lambda$ est nilpotente, donc il en est
de même de $s$, et le critère de (\hyperref[II.3.7.5-fr]{3.7.5}) s'applique.
\end{proof}

\begin{namedtheorem}[3.7.9]{Remarques}
\label{II.3.7.9-fr}
\begin{enumerate}
  \item[(i)] Avec les notations de (\hyperref[II.3.7.1-fr]{3.7.1}) et en
    tenant compte de (\hyperref[II.3.2.10-fr]{3.2.10}), on a, pour tout
    $n\in\mathbf{Z}$, un homomorphisme canonique
    \[
    \label{II.3.7.9.1-fr}
      \theta:r_{\mathcal{L},\psi}^*(\mathcal{O}_P(n))\longrightarrow
      \mathcal{L}^{\otimes n}
    \tag{3.7.9.1}
    \]
    défini de façon générale dans
    (\hyperref[II.3.5.6.2-fr]{3.5.6.2}). On voit aussitôt que sous les
    conditions de (\hyperref[II.3.7.3-fr]{3.7.3}), la restriction de cet
    homomorphisme à $X_{\psi^\flat(f)}$ s'explicite comme faisant
    correspondre à $s/f^k$ ($s\in S_{n+kd}$) l'élément
    $(\psi^\flat(s)|X_{\psi^\flat(f)})
    (\psi^\flat(f)|X_{\psi^\flat(f)})^{-k}$.
  \item[(ii)] Soit $\mathcal{F}$ un $\mathcal{O}_X$-Module quasi-cohérent, et
    supposons que $q$ soit quasi-compact et séparé, de sorte que pour tout
    $n\geq0$, $q_*(\mathcal{F}\otimes\mathcal{L}^{\otimes n})$ soit un
    $\mathcal{O}_Y$-Module quasi-cohérent
    (\hyperref[I.9.2.2-fr]{I, 9.2.2}). Soit
    $\mathcal{M}'=\bigoplus_{n\geq0}\mathcal{F}\otimes
    \mathcal{L}^{\otimes n}$, qui est un $\mathcal{S}'$-Module gradué
    quasi-cohérent, et considérons son image
    $\mathcal{M}=q_*(\mathcal{M}')=
    \bigoplus_{n\geq0}q_*(\mathcal{F}\otimes\mathcal{L}^{\otimes n})$
    (qui est un $\mathcal{S}$-Module quasi-cohérent, au moyen de
    l'homomorphisme $\psi^\flat$). Nous allons voir qu'il y a un
    homomorphisme canonique de $\mathcal{O}_X$-Modules
    \[
    \label{II.3.7.9.2-fr}
      \xi:r_{\mathcal{L},\psi}^*(\widetilde{\mathcal{M}})
      \longrightarrow\mathcal{F}|G(\psi).
    \tag{3.7.9.2}
    \]

    En effet, on a défini (\hyperref[II.3.5.6.1-fr]{3.5.6.1}) un
    homomorphisme canonique
    \[
    \label{II.3.7.9.3-fr}
      r_{\mathcal{L},\psi}^*(\widetilde{\mathcal{M}})
      \longrightarrow
      (q^*(\mathcal{M})\otimes_{q^*(\mathcal{S})}\mathcal{S}')^{\sim}
      |G(\psi),
    \tag{3.7.9.3}
    \]
    où le second membre est considéré comme un faisceau quasi-cohérent sur
    $\operatorname{Proj}(\mathcal{S}')$. D'autre part, on a un
    homomorphisme canonique
    \[
    \label{II.3.7.9.4-fr}
      q^*(q_*(\mathcal{M}'))\otimes_{q^*(\mathcal{S})}\mathcal{S}'
      \longrightarrow\mathcal{M}'
    \tag{3.7.9.4}
    \]
    pour tout $\mathcal{S}'$-Module gradué quasi-cohérent $\mathcal{M}'$~:
    pour tout ouvert $U$ de $X$, toute section $t'$ de
    $q^*(q_*(\mathcal{M}'_h))$ au-dessus de $U$ et toute section $b'$ de
    $\mathcal{S}'_k$ au-dessus de $U$, on fait correspondre à $t'\otimes b'$ la
    section $b'\sigma(t')$ de $\mathcal{M}'_{h+k}$, où $\sigma(t')$ est la
    section de $\mathcal{M}'_h$ au-dessus de $U$ qui correspond canoniquement
    (\hyperref[0.4.4.3-fr]{0, 4.4.3}) à $t'$. On en conclut un homomorphisme
    canonique
    \[
    \label{II.3.7.9.5-fr}
      (q^*(q_*(\mathcal{M}'))\otimes_{q^*(\mathcal{S})}\mathcal{S}')^{\sim}
      |G(\psi)\longrightarrow\widetilde{\mathcal{M}'}|G(\psi)
    \tag{3.7.9.5}
    \]
    et comme finalement $\widetilde{\mathcal{M}'}$ s'identifie
    canoniquement à $\mathcal{F}$
    (\hyperref[II.3.2.9-fr]{3.2.9, (i)}), on obtient bien l'homomorphisme
    canonique cherché.

    Sous les conditions de (\hyperref[II.3.7.3-fr]{3.7.3}), la restriction
    de cet homomorphisme à $X_{\psi^\flat(f)}$ s'explicite de la façon
    suivante~: étant donnée une section $t_{nd}$ de
    $\mathcal{F}\otimes\mathcal{L}^{\otimes nd}$ au-dessus de $X$, si
    $t'_{nd}$ est la section $t_{nd}$ considérée comme section de
    $q_*(\mathcal{F}\otimes\mathcal{L}^{\otimes nd})$ au-dessus de $Y$, à
    l'élément $t'_{nd}/f^n$ il correspond la section
    $(t_{nd}|X_{\psi^\flat(f)})
    (\psi^\flat(f)|X_{\psi^\flat(f)})^{-n}$ de $\mathcal{F}$ au-dessus de
    $X_{\psi^\flat(f)}$.
\end{enumerate}
\end{namedtheorem}

